\scanpage{226}
matrices \unsure{$\{$} $\underline u$ \unsure{$\in B$} such that $u\equiv1\ (2)$ i.e. $M\equiv0\ \unsure{(2)}\unsure{\}}$, by the \inserted{closed} subgroup $\underline L_0$ generated by $\ell_0=\begin{pmatrix}1&-2\\0&1\end{pmatrix}$, is identified with the group of matrices $\begin{pmatrix}\mu&0\\0&1\end{pmatrix}$ --- \inserted{\unsure{the}} maximal torus $T_0$ of $B$, more precisely we have a $\frac12$ direct product structure
\[
\begin{array}{ccccc}
\unsure{\underline B_1}&\simeq&T_0&\cdot&\underline L_0\\
\Vert&&\Vert&&\Vert\\
\left\{\begin{pmatrix}\mu&M\\0&1\end{pmatrix}\middle|\begin{array}{l}\mu\in\widehat{\mathbb Z}^{*}\\M\in2\widehat{\mathbb Z}\end{array}\right\}&&
\left\{\begin{pmatrix}\mu&0\\0&1\end{pmatrix}\middle|\mu\in\widehat{\mathbb Z}^{*}\right\}&&
\left\{\begin{array}{c}\ell_0^n\\\Vert\\\begin{pmatrix}1&-2n\\0&1\end{pmatrix}\end{array}\middle|n\in\widehat{\mathbb Z}\right\}
\end{array}
\tag{178}
\]

This shows that if we are interested in the quotient $\mathcal M_{0,3}^{!\sim}(0)/L_0\xrightarrow{\sim}\mathcal N_{0,3}^{!\sim}$, then the only abelian invariant we know how to associate with it in the present context (when we are not working with $\alpha,\beta$) is the multiplier $\mu$. But better still, we find \unsure{thanks} to the decomposition (178) in $\mathrm{GL}(2,\widehat{\mathbb Z})$ a canonical splitting of the extension $\mathcal M_{0,3}^{!\sim}(0)$ of $\mathcal N_{0,3}^{!\sim}$ by $L_0$ :

\uline{Proposition} For $\forall\,\dot u\in\mathcal N_{0,3}^{!\sim}$, there is a unique representative \unsure{$u=u_{\dot u}$|$u=u_{\xi}$} $\in\mathcal M_{0,3}^{!\sim}[0]$ such that we have $\underline u\in T_0$ i.e. $\underline u$ of the form $\begin{pmatrix}\mu&0\\0&1\end{pmatrix}$
\inserted{[i.e. $\xi=\begin{pmatrix}A&B\\C&D\end{pmatrix}$ of the form $\begin{pmatrix}A&0\\C&D\end{pmatrix}=\begin{pmatrix}\unsure{\ast|\frac{1}{D}}&0\\1-D&D\end{pmatrix}$]}, --- In other words, the subgroup

\margin{225}

\scanpage{227}
\mspage{413}
\deleted{\unsure{$\mathcal M_{0,3}^{!\sim}(0)$}}
\[
\mathcal M_{0,3}^{!\sim}(0)\overset{\mathrm{df}}=
\left\{u\in\mathcal M_{0,3}^{!\sim}(0)\ \middle|\ \underline u\in T_0\quad\text{i.e. }\underline u=\begin{pmatrix}\mu&0\\0&1\end{pmatrix}\right\}
\tag{179}
\]
is a section subgroup of the extension $\mathcal M_{0,3}^{!\sim}[0]$ of $\mathcal N_{0,3}^{!\sim}$ by $L_0$, hence also of the extension $\mathcal M_{0,3}^{!\sim}$ of $\mathcal N_{0,3}^{!\sim}$ by $\widehat\Pi_{0,3}$, \unsure{or} of the extension $\mathcal M_{0,3}^{\sim}$ of $\mathcal N_{0,3}^{!\sim}$ by $\widehat{\mathcal G}_{0,3}^{+}$.

\unsure{This|That} comes from the fact that if $u\in\mathcal M_{0,3}^{!\sim}[0]$ \inserted{defines} $\underline u=\begin{pmatrix}\mu&M\\0&1\end{pmatrix}$, thus with $M\equiv0\ (2)$, i.e. $M\in2\widehat{\mathbb Z}$, there is a \uline{unique} $n\in\widehat{\mathbb Z}$ such that $\underline{\ell_0^n u}=\begin{pmatrix}\mu&M-2n\\0&1\end{pmatrix}$ is in $T_0$ i.e. such that $\ell_0^n u\in\mathcal M_{0,3}^{!\sim}(0)$, namely $n=\dfrac M2$.

If $(\xi,\deleted{\unsure{\varepsilon}}\eta)$ satisfy conditions (133)(134) i.e. define a $u=u_{\xi,\eta}$ (cf. Theorem page~407'), then we shall say that $(\xi,\eta)$ is \uline{normalized} if we have $u_{\xi,\eta}\in\mathcal M_{0,3}^{!\sim}(0)$. Likewise, if an arbitrary element of $\mathcal M_{0,3}^{!\sim}$ is written in the form
\[
u=\cdot f^{-1}\cdot u_{\xi,\eta}=f^{-1}u_{\unsure{\xi}},
\]
we say that this representation \inserted{(\unsure{tainted} with the indeterminacy (136))} is \uline{normalized}, if $u_{\xi,\eta}=u_\xi$ is normalized, [i.e. \unsure{if} \inserted{for} $\xi\unsure{\simeq}\begin{pmatrix}A&B\\C&D\end{pmatrix}$ \deleted{the form} we have $B=0$]. Thus, \deleted{every} $u$ can be written uniquely in the form
\[
u_{f,\xi,\eta}=f^{-1}u_\xi,
\tag{\unsure{179}}
\]
\margin{226}

\scanpage{228}
with $\unsure{f|\xi}$ normalized. In fact, this corresponds to the $\frac12$-direct product decomposition
\[
\mathcal M_{0,3}^{!\sim}\simeq\mathcal M_{0,3}^{!\sim}(0)\cdot\widehat\Pi_{0,3},
\tag{180}
\]
moreover we likewise have
\[
\mathcal M_{0,3}^{\sim}\simeq\mathcal M_{0,3}^{!\sim}(0)\cdot\widehat{\mathcal G}_{0,3}^{+}.
\tag{181}
\]

\uline{Remark} We should at some point \unsure{try} to relate these canonical splittings to the choice of a uniformizer at $0$ of $\mathbb P^1_{\mathbb Q}=\widehat U_{0,3\mathbb Q}$. Of course, the same question also arises \unsure{at|for} the \inserted{other two} ``points at infinity'' $1,\infty$ of $U_{0,3\mathbb Q}$, relative to the section subgroups
\[
\mathcal M_{0,3}^{!\sim}(q)\qquad q\in\{0,1,\infty\}
\tag{182}
\]
corresponding to these points at infinity --- the construction of a section-group $\mathcal M_{0,3}^{!\sim}(0)$ in the normalizer of $L_0$, obviously applies without change for $L_1,L_\infty$. We shall moreover have
\[
\rho\bigl(\mathcal M_{0,3}^{!\sim}(q)\bigr)=\mathcal M_{0,3}^{!\sim}\bigl(\unsure{\dot\rho|\rho}(q)\bigr)\qquad\text{for }q\in\{0,1,\infty\},
\tag{183}
\]
so the \inserted{three} $\mathcal M_{0,3}^{!\sim}(q)$ are obtained from any one of them by applying $\rho,\rho^2$.
\margin{227}

\scanpage{229}
\mspage{415}
47. Return to the homomorphism $\widehat\Pi_{0,3}\xrightarrow{\widehat\varphi}\widehat{\mathcal G}_{0,3}^{+}$ and the condition of commutation with \uline{it} up to inner automorphisms. The groups $\mathcal M_{0,3}^{!\approx},\mathcal M_{0,3}^{\approx}$ and the homomorphism $\Phi:\mathcal M_{0,3}^{!\approx}\longrightarrow\mathcal M_{0,3}^{\approx}$.

Recall \inserted{(\S41)} that $\widehat\varphi$ is described by
\[
\widehat\varphi(\ell_0)=\varepsilon_0=\sigma_\infty\rho,\quad
\widehat\varphi(\ell_1)=\sigma_\infty,\quad
\widehat\varphi(\ell_\infty)=\rho^{-1}.
\tag{1}
\]
Let
\[
u=u_{\alpha,\beta,f}\quad\in\mathcal M_{0,3}^{!\sim}
\]
be a loop automorphism of $\widehat\Pi_{0,3}$, commuting with $\mathcal G_3$ up to inner automorphisms, defined by
\[
\begin{aligned}
u(\sigma_\infty)&=h\sigma_\infty=\operatorname{int}(\beta_0)\sigma_\infty
&\left(\begin{array}{l}h=\beta_0\sigma_\infty(\beta_0)^{-1},\\\beta_0=f^{-1}\beta\end{array}\right)
\end{aligned}
\tag{2}
\]
\[
\begin{aligned}
u(\rho)&=g\rho=\operatorname{int}(\alpha_0)\rho
&\left(\begin{array}{l}g=\alpha_0\rho(\alpha_0)^{-1},\ \deleted{\unsure{=}}\\\alpha_0=f^{-1}\alpha\end{array}\right)
\end{aligned}
\tag{3}
\]
and let us write that $u$ commutes with $\varphi$ up to inner automorphisms, i.e. that there exists \deleted{\unsure{$\lambda\in$}}
\[
\lambda\in\widehat\Pi_{0,3}
\tag{4}
\]
such that we have
\[
\widehat\varphi u(x)=\operatorname{int}(\lambda)u\widehat\varphi(x)
\qquad\forall x\in\widehat\Pi_{0,3}.
\tag{5}
\]
It suffices to check (5) for $x=\ell_0$ and $x=\ell_1$.

1) Case $x=\ell_0$ \deleted{\unsure{$=(\sigma_\infty\rho)^2$}}, whence
\[
u(x)=\deleted{u(\sigma_\infty)u(\rho)}\deleted{\unsure{=h\sigma_\infty g\rho=h(\sigma_\infty g\sigma_\infty^{-1})\sigma_\infty\rho}}
\inserted{\operatorname{int}(f^{-1})(\ell_0^\mu)}
\]
\[
\widehat\varphi(u(x))=\operatorname{int}\bigl(\widehat\varphi(f)^{-1}\bigr)\widehat\varphi(\ell_0^\mu)=\operatorname{int}\bigl(\widehat\varphi(f)^{-1}\bigr)\varepsilon_0^\mu
\]
\[
u\widehat\varphi(x)=u(\varepsilon_0)=\operatorname{int}(f^{-1})(\varepsilon_0^\mu)
\]
so the relation (5) reads
\[
\operatorname{int}\bigl(\deleted{\unsure{f\lambda^{-1}}}\widehat\varphi(f)\unsure{\lambda}f^{-1}\bigr)(\varepsilon_0^\mu)=\varepsilon_0^\mu
\]
\margin{228}

\scanpage{230}
is equivalent to
\[
\deleted{\unsure{f\lambda^{-1}}}\widehat\varphi(f)\lambda f^{-1}\in L_0^{\sim}\quad
\left(=\operatorname{Centr}_{\widehat\Pi_{0,3}}(\varepsilon_0^\mu)\right)
\]
\margin{N.B. we do not necessarily have $\widehat\varphi(f)\in\widehat\Pi_{0,3}$, so \unsure{neither} \illegible{1} $\widehat\varphi(f)\lambda f^{-1}\in\widehat\Pi_{0,3}\ldots$}
i.e. to the existence of $r\in\widehat{\mathbb Z}$ such that we have
\[
\deleted{\unsure{f\lambda^{-1}}}\widehat\varphi(f)\lambda f^{-1}=\varepsilon_0^r,
\]
i.e.
\[
\lambda=\widehat\varphi(f)^{-1}\varepsilon_0^r f.
\tag{6}
\]
This thus makes it possible to eliminate $\unsure{\exists\lambda}$ from the second relation :

2) Case $x=\ell_1$, whence
\[
u(x)=\operatorname{int}(f_1^{-1})(\ell_1^\mu)
\qquad\bigl(f_1^{-1}=\deleted{\unsure{\rho(f)}}g\rho(f^{-1})\bigr)
\]
\[
\widehat\varphi(u(x))=\operatorname{int}\bigl(\widehat\varphi(f_1)^{-1}\bigr)\sigma_\infty^\mu=\operatorname{int}\bigl(\widehat\varphi(f_1)\bigr)\sigma_\infty
\]
(taking into account that $\sigma_\infty^\mu=\sigma_\infty$ since $\mu\equiv1\ (2)$).
\deleted{\unsure{$u\widehat\varphi(x)=\operatorname{int}(\widehat\varphi(f_1))$}}
\[
u\widehat\varphi(x)=u(\sigma_\infty)=\operatorname{int}(\beta_0)(\sigma_\infty),
\]
So the relation (5) reads
\[
\operatorname{int}\bigl(\widehat\varphi(f_1)^{+1}\lambda\beta_0\bigr)\sigma_\infty=\sigma_\infty,
\]
which also means that there exists $\unsure{i|\alpha}\in\mathbb Z/2\mathbb Z$ such that we have \deleted{\illegible{1}}
\[
\widehat\varphi(f_1)^{+1}\lambda\beta_0=\sigma_\infty^i
\quad\text{i.e.}\quad
\widehat\varphi(f_1^{-1})^{+1}=\lambda\beta_0\sigma_\infty^i,
\]
and replacing $\lambda$ by its value (6), and $f_1^{-1}$ by its value $g\rho(f^{-1})\deleted{\unsure{f}}$, we find the condition
\[
\deleted{\unsure{\widehat\varphi(\rho(f)^{-1})\widehat\varphi(g)^{-1}=\widehat\varphi(f)^{-1}\ell_0^r f\beta_0\sigma_\infty^\alpha}}
\]
\[
\widehat\varphi\bigl(g\rho(f)^{-1}\bigr)=\widehat\varphi(f)^{-1}\varepsilon_0^r\underbrace{f\beta_0}_{\beta}\sigma_\infty^{\inserted{\unsure{\alpha}}}\quad\text{i.e.}
\]
\[
\widehat\varphi\bigl(\underbrace{fg\rho(f)^{-1}}_{\xi=\alpha\rho(\alpha)^{-1}}\bigr)=\varepsilon_0^r\beta\sigma_\infty^\alpha
\]
\margin{229}

\scanpage{231}
\mspage{416}

Thus :

\uline{\deleted{Proposition}\inserted{Theorem}} For
\[
\widehat\varphi:\widehat\Pi_{0,3}\longrightarrow
\widehat{\unsure{\mathcal G}}_{0,3}^{+}
\]
to commute up to inner automorphisms with the element
$u\in\mathcal M_{0,3}^{!\sim}$, corresponding to the invariants
$g,h,f,\alpha_0,\beta_0,\alpha,\beta,\xi,\eta$ (cf. \S 46),
it is necessary and sufficient that we have
\[
\widehat\varphi(\xi)=\varepsilon_0^r\beta\sigma_\infty^{\unsure{i|\alpha}}
\qquad\text{for }r\in\widehat{\mathbb Z},\quad
\alpha\in\mathbb Z/2\mathbb Z\ \text{suitable}.
\tag{7}
\]
\margin{\uline{N.B.} $r,\alpha$ are well determined when $\alpha$ and $\beta$
(or again $\xi$ and $\beta$) are chosen.}
We shall then have (5), with \deleted{Note that}
$\lambda$ given by (6) ;
$\lambda=\widehat\varphi(f)^{-1}\ell_0^r f$.

Note that $u$ defines the triple $(\alpha,\beta,f)$ only modulo
a transformation of the form
\deleted{$(\alpha,\beta,f)\ \unsure{\mapsto}\ \unsure{\alpha'=\ell_0^n\alpha,\ \beta}$}
\[
(\alpha,\beta,f)\longmapsto
(\alpha'=\ell_0^n\alpha,\ \beta'=\ell_0^n\beta,\ f'=\ell_0^n f)
\tag{8}
\]
where $n\in\widehat{\mathbb Z}$. We shall then have (putting
$\xi'=\alpha'\rho(\alpha')^{-1}=\ell_0^n\xi\ell_1^{-n}$)
\[
\widehat\varphi(\xi')=\varepsilon_0^{r'}\beta'
\sigma_\infty^{\unsure{i'|\alpha'}}
\tag{*}
\]
with $r'\in\widehat{\mathbb Z}$,
$\unsure{i'|\alpha'}\in\mathbb Z/2\mathbb Z$ also well determined.
Now we shall have
\[
\begin{aligned}
\widehat\varphi(\unsure{\xi'})
 &=\widehat\varphi(\ell_0^n\unsure{\xi}\ell_1^{-n})\\
 &=\widehat\varphi(\ell_0^n)\widehat\varphi(\unsure{\xi})
   \widehat\varphi(\ell_1^{-n})\\
 &=\varepsilon_0^n\widehat\varphi(\unsure{\xi})\sigma_\infty^{-n}
\end{aligned}
\]
\deleted{\illegible{1}}\deleted{\illegible{1}}
and the relation (*) reads
\[
\widehat\varphi(\xi)=\varepsilon_0^{\unsure{r'-n|r'-2n}}\beta'
\sigma_\infty^{\unsure{i'-n|\alpha'-n}}
\]
or again, since $\beta'=\ell_0^n\beta=\varepsilon_0^{2n}\beta$
\[
\widehat\varphi(\xi)=\varepsilon_0^{r'+n}\beta
\sigma_\infty^{\unsure{\alpha'-n|i'-n}},
\]
which, compared with (7), gives
\[
r'=r-n,\qquad\unsure{i'|\alpha'}=\unsure{i|\alpha}+n\ \mathrm{mod}\,2.
\tag{9}
\]
\margin{230}